\chapter{いつ書き、いつ読むか：\nt{ACET}を加える}
\chaptermark{\nt{ACET}を加える}
\label{sec:acet}

\section{足りないもの}

メモリチップには、アドレス線とデータ線のほかにもう1本の線、書き込み許可がある。これがオフの間、メモリは読み出されるだけである。オンになると、データ線上にあるものが書き込まれる。この線がなければメモリは働かない。読み出すたびに同時に何かが書き込まれるなら、読んだもの、誤って読んだものまで、すぐに書き戻されてしまう。

脳ではこの問題はチップよりも深刻である。第\ref{sec:glutcomputer}章では、記憶とは自身の結合によってオンに保たれる\nt{GLUT}ニューロンの群だった（図~\ref{fig:bistable}）。第\ref{sec:dofa}章では、これらの結合が動作中にそのままヘッブ則で学習することを見た。つまり同じシナプスが古いものを保持し、新しいものも受け入れる。そして独立した書き込みフェーズはなく、クロックもない。ネットワークは、見ているものを記憶すべき瞬間と、知っていることを思い出すべき瞬間をどう区別するのか。第\ref{sec:neurontypes}章で、この線を担うのは\nt{ACET}だと仮定した。この章では、そのような線がどう働くべきか、なぜそれなしでは済まないのかを確かめる。

\section{1本の配線の三つの作用}

脳の内部で働くアセチルコリンニューロンは少なく、いくつかの小さな群にまとまっているが、その出力は皮質全体に広がる。これは\nt{DOPA}や\nt{NORA}と同じブロードキャストチャネルである。皮質のアセチルコリンレベルは注意深い覚醒時に高く、徐波睡眠で低い~\cite{Hasselmo1999}。ドーパミンと同じく、信号の意味は受け手の受容体が決める（命題~\ref{prop:receptor}）。アセチルコリンには、チャネルを開く速い受容体と、細胞内の反応を起動する遅い受容体の両方がある。ここで重要な作用は三つある。

\emph{第一に、内部結合を弱める。}嗅皮質のスライスでの実験で、アセチルコリンは同じ領域のニューロン間の結合による伝達を強く抑え、外から入ってくる入力にはほとんど触れないことが示された~\cite{HasselmoBower1992}。

\emph{第二に、入力を強める。}アセチルコリンはニューロンの興奮性を高め、いくつかの入力経路の伝達を促進する。感覚器官からの信号がネットワーク自身の活動より大きくなる~\cite{Hasselmo2006}。

\emph{第三に、重みの変化を容易にする。}アセチルコリンのレベルが高いと、結合は変わりやすくなる~\cite{Hasselmo2006}。

エンジニアにとって、三つの作用はすべて一つの操作、「読む」モードから「書く」モードへの切り替えである。ネットワークは自分自身を聞くのをやめ、入力を聞き始め、聞いたことを記憶する。

\section{補完するネットワーク}

互いに結合した$N$個の\nt{GLUT}ニューロンをとる。その結合に、いくつかのパターン（同時に活動する$pN$個のニューロンの組）をヘッブ則で保存する~\cite{Hebb1949}。このネットワークにパターンの半分を与えると、結合が残りの半分を興奮させる。ネットワークは、図~\ref{fig:bistable}の群が自身を維持したのと同じように、一部からパターンを\emph{補完する}。これは連想記憶であり、内容の一部がアドレスになる。第\ref{sec:gaba}章と同じく、\nt{GABA}は活動するニューロンの総数を$pN$程度に保ち、二つのパターンが同時に点灯しないようにする。

アセチルコリンのレベルを$0$から$1$の$a$で表し、その三つの作用をそのままモデルに書き込む。
\begin{empheq}[box=\keybox]{equation}
\tau\frac{\dd r_i}{\dd t} = -r_i + F\Bigl(\tfrac{1+a}{2}\,x_i + (1-a)\sum_j W_{ij}\,r_j - g\Bigr),
\qquad
\frac{\dd W_{ij}}{\dd t} = \eta\,a\,(r_i-p)(r_j-p) .
\label{eq:acetnet}
\end{empheq}
ここで$r_i$はニューロン$i$の発火率、$x_i$は感覚器官からのその入力、$F$は閾値型の伝達特性、$g$は\nt{GABA}による全体的な抑制で、活動ニューロンが$pN$より多いと増える。係数$\frac{1+a}{2}$は入力の強化、$1-a$は内部結合の弱化、$\eta a$は学習率である。第二式はまばらなパターンに便利な形のヘッブ則~\eqref{eq:hebb}である~\cite{Tsodyks1988}。両方のニューロンがふだんより活動すると重みは増え、一方だけが活動すると減る。重みの負の部分は、いつものように\nt{GABA}仲介ニューロンが担う。クロックはない。ニューロンも重みも連続的に変わる。

\begin{figure}[t]
\centering
\begin{tikzpicture}[>=Stealth,
  nrn/.style={circle,draw,very thick,fill=blue!6,minimum size=0.8cm,inner sep=0pt,font=\small},
  acet/.style={circle,draw,very thick,double,double distance=1.2pt,fill=orange!15,minimum size=1.35cm,inner sep=0pt,font=\scriptsize},
  bc/.style={->,thick,densely dashed,orange!85!black}]
\foreach \i/\y in {1/2.4,2/1.2,3/0}{
  \node[nrn] (N\i) at (3,\y) {$r_\i$};
  \draw[->,thick,blue!60!black] (0,\y) node[left,black] {\small $x_\i$} -- (N\i);}
\node[left,align=right,font=\small] at (-0.5,1.2) {感覚器官\\から};
\draw[->,thick,gray!60!black] (N1) to[bend left=30] (N2);
\draw[->,thick,gray!60!black] (N2) to[bend left=30] (N1);
\draw[->,thick,gray!60!black] (N2) to[bend left=30] (N3);
\draw[->,thick,gray!60!black] (N3) to[bend left=30] (N2);
\draw[->,thick,gray!60!black] (N1.east) .. controls (4.6,2.0) and (4.6,0.4) .. (N3.east);
\node[right,font=\small,gray!40!black,align=left] at (4.7,1.2) {内部\\結合 $W$};
\node[acet] (A) at (1.2,-1.9) {\nt{ACET}};
\draw[bc] (A) -- (1.6,-0.15);
\node[font=\small,orange!60!black,anchor=east] at (0.95,-0.9) {$+$ 入力};
\draw[bc] (A) -- (4.3,0.6);
\node[font=\small,orange!60!black,anchor=west] at (4.3,0.1) {$-$ 内部結合、$+$ 学習率};
\node[right,font=\small,align=left] at (2.2,-1.9) {高レベル：「書く」\\低レベル：「読む」};
\end{tikzpicture}
\caption{書き込み許可線としての\nt{ACET}。ニューロン$r_i$は感覚器官からの入力$x_i$を受け（青い矢印）、パターンを保存している内部結合$W$（灰色）を通じて互いに興奮させる。アセチルコリン（破線の矢印）は、\eqref{eq:acetnet}のように入力を強め、内部結合を弱め、重みの変化を容易にする。}
\label{fig:acetnet}
\end{figure}

\section{書き込みと読み出しは互いに邪魔をする}

書き込みと読み出しが本当に衝突する課題で、モデル~\eqref{eq:acetnet}を確かめよう。$200$個のニューロンのネットワークが、活動ニューロン$20$個のパターンを五つすでに保存している。ここで六つ目を記憶しなければならないが、それは古いパターンの一つと半分重なっている。10個のニューロンが共通である。これはいつも起きることである。新しい部屋はなじみの部屋に似ていて、新しい顔は古い顔に似ている。

\emph{低い$a$での書き込み。}まずアセチルコリンの最初の二つの作用を第三の作用から切り離す。学習率はフルのままにし、変えるのは入力の強化と内部結合の弱化だけにする。$a=0$では、入力が共通の10個のニューロンを点灯させ、それらが内部結合を通じて\emph{古い}パターンの残り半分を点灯させる。\nt{GABA}は両方が同時に点灯することを許さず、自身の結合に支えられた古いパターンが、入力だけに支えられた新しいパターンを押しのける。ネットワークは目の前にあるものではなく覚えているものを見て、ヘッブ則はその見たものを律儀に書き込む。

結局、新しいパターンはまったく記憶されなかった。その特有の半分を与えても、ネットワークは何も思い出さない。そして古いパターンは壊れている。特有の半分で呼び出すと、平均して10個中6個の無関係なニューロンを引き連れてくる。$a=0.1$では新しいパターンは記憶されるようになるが、古いものと融合し、破損はさらにひどくなる。二つのそれぞれが、自分の半分で呼び出されると約8個の無関係なニューロンを引き連れる。$a=0.2$からは書き込みはきれいになり、想起時の余分なニューロンは1個未満になる（10回の実行の平均）。

\emph{本来の学習率での書き込み。}今度はアセチルコリンが、\eqref{eq:acetnet}のように学習率も決めるとする。1回の提示で新しいパターンが丸ごと記憶されるのは$a\ge0.9$の場合だけである。$a=0.75$では10分の8、$a\le0.5$ではまったく記憶されない。

\emph{読み出し。}五つのパターンがきれいに書き込まれたネットワークにパターンの半分を与え、それから取り去る。$a\le0.2$ではネットワークはパターンを丸ごと補完し、自力で保つ。$a=0.3$ではほぼ丸ごと。$a\ge0.4$では内部結合が弱められすぎていて、手がかりが消えると活動も消える（図~\ref{fig:acetsim}）。

\begin{figure}[t]
\centering
\begin{tikzpicture}[>=Stealth,x=9cm,y=3.2cm]
\draw[->,gray!70!black] (0,0) -- (1.1,0) node[right,black] {\small $a$};
\draw[->,gray!70!black] (0,0) -- (0,1.2);
\foreach \x in {0,0.2,0.4,0.6,0.8,1.0}{\draw[gray!60!black] (\x,-0.02) -- (\x,0.02); \node[below,font=\scriptsize] at (\x,0) {$\x$};}
\foreach \y in {0,0.5,1}{\draw[gray!60!black] (-0.01,\y) -- (0.01,\y); \node[left,font=\scriptsize] at (0,\y) {$\y$};}
\node[rotate=90,font=\small] at (-0.1,0.6) {復元されたパターンの割合};
\draw[very thick,blue!60!black] plot[mark=*,mark size=1.6pt] coordinates {(0,1.00) (0.1,1.00) (0.2,1.00) (0.3,0.96) (0.4,0.04) (0.5,0.00) (0.75,0.00) (1.0,0.00)};
\draw[very thick,red!70!black] plot[mark=*,mark size=1.6pt] coordinates {(0.1,0.00) (0.2,0.00) (0.3,0.00) (0.5,0.00) (0.75,0.80) (0.9,1.00) (1.0,1.00)};
\node[font=\small,blue!60!black,anchor=west,align=left] at (0.03,0.5) {読み出し：\\半分から\\パターンを};
\node[font=\small,red!70!black,anchor=east,align=right] at (1.0,0.35) {書き込み：\\新パターンを\\1回で};
\end{tikzpicture}
\caption{モデル~\eqref{eq:acetnet}：ニューロン$200$個、活動ニューロン$20$個のパターン、10回の実行の平均。青い点は読み出し。手がかりを取り去った後、半分からパターンのどれだけの割合が復元されたか。赤い点は書き込み。アセチルコリンが学習率も決める場合に、古いパターンと半分似た新しいパターンのどれだけの割合が1回の提示の後に思い出されるか。読み出しは$a\le0.3$で、書き込みは$a\ge0.9$で働く。}
\label{fig:acetsim}
\end{figure}

\begin{keyblock}
\begin{proposition}[書き込みと読み出しには別々のモードが要る]
\label{prop:writeread}
同じ結合が古いものを保持し新しいものを学習し、アセチルコリンが学習率も決めるモデル~\eqref{eq:acetnet}では、書き込みと読み出しがどちらも同じようにうまくいくレベル$a$は存在しない。書き込みには内部結合を弱めなければならない。さもないとネットワークは入力ではなく思い出したものを書き込む。読み出しには内部結合を強めなければならない。さもないと一部からパターンを補完できない。スイッチが必要であり、1本のブロードキャスト配線がその役を果たしうる。
\end{proposition}
\end{keyblock}

アセチルコリンが学習率を変えないなら、両方の仕事に使える狭い帯$0.2\le a\le0.3$がある。しかしそのときネットワークは読み出しのたびにも学習し、読んだものをその誤りごと書き直してしまう。学習中に内部結合を抑えるという考え自体は、スライスでの測定に基づくモデルからとったものである~\cite{HasselmoAndersonBower1992}。上の数値は私たちの簡略化したモデルについてのものであり、脳でモードの境界が正確にどこにあるかは分かっていない。

\section{目をどこまで信じるか}

「書く／読む」のスイッチは、より一般的な問い、入力をどれだけ信じ、記憶をどれだけ信じるかの極端な場合である。この問いにはエンジニアの厳密な答えがあり、それは1本の配線がなぜモードと学習率を同時に制御するのかを明らかにする。

ネットワークが、ゆっくりさまよう量$s(t)$を追跡しているとする。その分散は1秒あたり$q$だけ増える。感覚器官は$y(t)=s(t)+$（強度$R$のノイズ）を報告する。連続時間での最良の推定$\hat s$はカルマン--ビューシーフィルタで与えられる~\cite{KalmanBucy1961}。
\begin{empheq}[box=\keybox]{equation}
\frac{\dd\hat s}{\dd t} \;=\; \frac{P}{R}\,\bigl(y-\hat s\bigr),
\qquad
\frac{\dd P}{\dd t} \;=\; q-\frac{P^2}{R} ,
\label{eq:kalman}
\end{empheq}
ここで$P$は推定誤差の分散、つまりネットワーク自身が知っていることにどれだけ自信がないかである。第一式はニューロンモデル~\eqref{eq:glutneuron}の膜と同じRC回路だが、時定数$R/P$が変わる。ネットワークに自信がないと$P$は大きく、時定数は小さく、推定はすばやく入力に従う。これが「書く」である。自信があると$P$は小さく、推定はほとんど入力を聞かずに記憶にしがみつく。これが「読む」である。

定常状態では$P=\sqrt{qR}$で、記憶の時定数は$\sqrt{R/q}$である。世界が100秒で1単位ずれ、$q=0.01$、感覚器官のノイズが$R=1$なら、記憶は約10秒保つべきである。

ここで重要なのは、一つの数$P/R$が二つのことを同時に決めることである。記憶に対する入力の重みと、保持している推定が変わる速さである。これはまさに\eqref{eq:acetnet}でアセチルコリンがすることである。仮説~\cite{YuDayan2005}によれば、アセチルコリンはまさに$P$に似た量、\emph{予期された}不確実性、つまりネットワークがあらかじめ知っている不確実性をネットワークに伝える。このチャネルが常に遅いわけではない。\nt{ACET}ニューロンの一部は報酬と罰に20ミリ秒ほどで応答し、強化が予期しないものであるほど強く応答する~\cite{Hangya2015}。

\begin{keyblock}
\begin{proposition}[入力の重みと学習率に一つのつまみ]
\label{prop:kalman}
最適フィルタ~\eqref{eq:kalman}では、入力を記憶と混ぜる重みと学習率は同じ量$P/R$である。したがってそれらを一つの信号で制御するのは合理的である。世界の性質が変わらなければ、この信号は一定レベル$\sqrt{q/R}$に落ち着き、調整が必要なのは世界の性質が変わったときだけである。
\end{proposition}
\end{keyblock}

命題の第二文は第\ref{sec:nora}章の結果を説明する。そこでは驚きによる学習率の調整は最良の一定値に勝てなかった。変化の性質が一定の世界では、最良の学習率は一定値そのものである。利得が生じるのは、世界が突然別物になったときである。そのとき推定は急に入力から大きく外れるが、$P$はそれを知らない。正しい行動は$P$を急に上げ、それによって書き込みモードに移ることである。第\ref{sec:nora}章で論じた仮説によれば、そのような\emph{予期しない}変化を伝えるのは\nt{NORA}である~\cite{YuDayan2005}。分業は自然に決まる。\nt{NORA}は世界が変わったことに気づき、\nt{ACET}は新しい状況が学習されるまでネットワークを書き込みモードに保つ。

\section{読み出しモードとしての睡眠}

アセチルコリンの高レベルが書き込みモードなら、低レベルのときネットワークは何をするのか。徐波睡眠ではアセチルコリンのレベルが下がり、内部結合は抑制から解放され、感覚器官からの入力はほとんどない。読み出しモードのネットワークは、手がかりなしに、自分に書き込まれているものを再生する。活動の記録によれば、日中一緒に働いたニューロンは徐波睡眠でふたたび一緒に発火し~\cite{WilsonMcNaughton1994}、しかも同じ順序で発火する~\cite{Skaggs1996}。仮説~\cite{Hasselmo1999}によれば、これらの再生は日中にすばやく学習したものを長期記憶に書き写す（再生の短いバーストを人工的に延ばすと記憶がよくなる~\cite{FernandezRuiz2019}）。日中は入力からの書き込み、夜間は内部からの書き直しである。エンジニアにとっては、日中に高速バッファに書き込み、夜間にそれを低速の不揮発性メモリにフラッシュするのに似ている。

\section{まとめ}

\begin{table}[t]
\centering
\footnotesize
\setlength{\tabcolsep}{4pt}
\renewcommand{\arraystretch}{1.15}
\begin{tabular}{>{\raggedright}p{3.2cm}>{\raggedright}p{4.2cm}>{\raggedright\arraybackslash}p{6.0cm}}
\toprule
\nt{ACET}がすること & 方法 & 分かっていること \\
\midrule
内部結合を弱める & \eqref{eq:acetnet}の係数$1-a$ & スライスで測定済み \\
入力を強める & 係数$\frac{1+a}{2}$ & 興奮性と入力経路の伝達の増大は測定済み \\
学習を容易にする & 学習率$\eta a$ & 測定済み \\
「書く／読む」を切り替える & 命題~\ref{prop:writeread} & モデルから導かれる。モードの直接測定はない \\
予期された不確実性を伝える & \eqref{eq:kalman}の$P$ & 仮説 \\
睡眠中の再生のためにネットワークを解放する & 徐波睡眠での低い$a$ & 睡眠中の低レベルと再生は測定済み。長期記憶への書き直しは仮説 \\
\bottomrule
\end{tabular}
\caption{\nt{GLUT}、\nt{GABA}、\nt{DOPA}、\nt{NORA}のネットワークに\nt{ACET}が加えるもの。}
\label{tab:acet}
\end{table}

\nt{DOPA}はネットワークに重みをどこへ変えるかを伝え、\nt{NORA}は知っていることをどれだけ断固として使うかを伝え、\nt{ACET}は世界を聞くか自分を聞くかを伝える（表~\ref{tab:acet}）。これは表~\ref{tab:types}の最後のブロードキャスト制御線、書き込み許可である。これがなければ、パターンを読み出すのと同じ結合にパターンを保存する記憶は、読み出すたびに自分を壊してしまう。

\section{問題}

\begin{exercise}[\lvlA]
\label{ex:09:1}
\emph{どれだけ覚えておくか。}$q=0.01$、$R=1$のフィルタ~\eqref{eq:kalman}は約$10$~s覚えている。(a)~センサノイズの振幅が2倍になった場合、(b)~世界のさまよいが100倍速くなった場合について、$P_\infty$と記憶$\sqrt{R/q}$を求めよ。(c)~ネットワークが1分覚えているには、ノイズは何倍にならなければならないか。
\end{exercise}

\begin{exercise}[\lvlB]
\label{ex:09:2}
\emph{変化の後の書き込みモード。}世界が変わり、$q=0.01$、$R=1$のフィルタ~\eqref{eq:kalman}の不確実性が$P_0\to\infty$に跳ね上がった。(a)~$P$はどんな時定数で$P_\infty$に戻るか。(b)~何秒後に学習率が定常値を$10\,\%$以内でしか上回らなくなるか。
\end{exercise}

\begin{exercise}[\lvlB]
\label{ex:09:3}
\emph{一定の学習率。}ネットワークは$\dd\hat s/\dd t=(y-\hat s)/T$で学習する。世界はさまよい、$\dd s/\dd t=w$、$y=s+n$で、$w$と$n$は強度$q$と$R$の白色ノイズである。最良の$T$とそのときの誤差の分散を求めよ。$T$が$2$倍、$10$倍ずれていると、分散は何倍になるか。
\end{exercise}

\begin{exercise}[\lvlB]
\label{ex:09:4}
\emph{書き込みの境界はどこか。}\texttt{models/\allowbreak acet\_memory.py}では閾値$\theta=0.35$、パターン内の重み$w=0.041$（理想的には$0.045$）である。新しいパターンは古いパターンと$m=10$個のニューロンを共有する。\eqref{eq:acetnet}でどの$a$なら、古いパターンの残りのニューロンがこの$m$個から点灯しないか（閾値は鋭いとする）。
\end{exercise}

\begin{exercise}[\lvlB]
\label{ex:09:5}
\emph{シミュレーション：記憶容量。}\texttt{models/\allowbreak acet\_memory.py}の関数\texttt{read\_sweep}を変更せよ。$a=1$で$M$個のパターン（$M$は$5$から$100$）を書き込み、次に$a=0$で最初の10個を半分から読み出す。どの$M$で誤りが現れるか。限界$M_{\max}$は$N$と$p$にどう依存するか。
\end{exercise}

\begin{exercise}[\lvlB]
\label{ex:09:6}
\emph{設計：漏れのない書き込み。}学習率$\eta a$では、ネットワークは読み出し時にも学習する。$a\le0.3$でゼロ、$a\ge0.9$で$\eta a$に劣らない単調な$\eta(a)\le\eta$を提案せよ。検証：無関係なニューロンを3個含む手がかりで$a=0.2$で$20$回読み出した後、そのうち何個がパターンに入ったか。
\end{exercise}

\begin{exercise}[\lvlC]
\label{ex:09:7}
\emph{夜間にバッファをどう書き直すか。}(a)~睡眠中は$a=0.05$で、再生は$0.1\,\text{s}$続く（日中は$a=1$で$0.2\,\text{s}$）。学習率$\eta a$と問題~\ref{ex:09:6}の$\eta(a)$のそれぞれで、日中の1回の提示と同じ重み変化を蓄積するには何回の再生が要るか。(b)~ストレージはバッファより$100$倍遅く学習し、一晩にパターンを$0.1\,\text{s}$ずつ$20$回聞く。何晩必要か。
\end{exercise}
